%!TEX encoding = IsoLatin

%
% Chapitre "Introduction"
%

\chapter{Introduction}
\label{s:intro}
Le suivi de la faune arctique représente un défi important en raison des conditions climatiques extrêmes, telles que les tempêtes de neige fréquentes et les variations de température, auxquelles s’ajoutent l’isolement géographique et l’accès limité au territoire. Le ministère de la faune arctique souhaite améliorer la fiabilité des données de suivi concernant les populations de petits mammifères, notamment les lemmings, dont l’activité sous la neige constitue un indicateur essentiel de l’état des écosystèmes du Grand Nord. C’est pourquoi celui-ci lance le projet de développement d’un concept de capteur autonome fixe pour le comptage et la documentation de la faune arctique. \\

L’équipe NordiQc a été contactée avec le mandat de produire le design conceptuel du système autonome permettant grossièrement de détecter et d’enregistrer les événements liés à l’activité des lemmings, de recueillir des vidéos afin de collecter des informations à des fins statistiques, d’archiver les données et d’assurer l’interaction minimale externe avec le système en cas de problème. \\

Ce rapport est structuré de manière à présenter d’abord les besoins du client, puis à établir les objectifs en lien avec ceux-ci et le cahier des charges. Il expose ensuite les trois concepts proposés et leur analyse de faisabilité, suivis d’une analyse quantitative des concepts globaux. Enfin, il se conclut par une description détaillée du concept retenu à l’issue des évaluations de l’équipe contactée. 




